% Part: turing-machines
% Chapter: machines-computations
% Section: halting-states

\documentclass[../../../include/open-logic-section]{subfiles}

\begin{document}

\olfileid{tur}{mac}{hal}
\olsection{Kahanan Mandheg}

\begin{explain}
Sanajan kita wis netepake mesin mung mandheg
nalika ora ana instruksi kang bisa ditindakake,
gambaran umum mesin Turing nduwé
\emph{kahanan mandheg}~$h$ mligi,
saengga $h \in Q$.

Gagasan kahanan mandheg iku prasaja: nalika mesin
wis ngrampungake pakaryane (wis siyap mutusake input iku ditampa,
utawa wis rampung nulis output), mesin mlebu
kahanan~$h$ lan mandheg. Sawetara mesin nduwé
rong kahanan mandheg, siji kanggo nampa input
lan siji kanggo nolak input.
\end{explain}

\begin{ex}\emph{Kahanan Mandheg}.
Kanggo njlentrehake gagasan iki, ayo diwiwiti
saka owah-owahan ing mesin genep. Tinimbang
mesin mandheg ing kahanan~$q_0$ nalika input
genep, kita bisa nambah instruksi kanggo
ngirim mesin menyang kahanan mandheg.
\[
\begin{tikzpicture}[->,>=stealth',shorten >=1pt,auto,node distance=2.8cm,
                    semithick]
  \tikzstyle{every state}=[fill=none,draw=black,text=black]

  \node[initial, state]         (A)                     {$q_0$};
  \node[state]         (B) [right of=A] {$q_1$};
  \node[state]         (C) [below of=A] {$h$};

  \path (A) edge [bend left] node {\TMtrans{\TMstroke}{\TMstroke}{R}} (B)
  	    edge node {\TMtrans{\TMblank}{\TMblank}{N}} (C)
        (B) edge [loop above] node {\TMtrans{\TMblank}{\TMblank}{R}} (B)
            edge [bend left] node {\TMtrans{\TMstroke}{\TMstroke}{R}} (A);
\end{tikzpicture}
\]

Ayo tuladha iki dikembangake maneh. Nalika mesin
nemtokake yen input ganjil, mesin ora tau mandheg.
Kita bisa ngowahi mesin kanthi nambah
\emph{kahanan nolak}, yaiku ngganti instruksi
kang bola-bali mau karo instruksi kanggo
mlebu kahanan nolak~$r$.
\[
\begin{tikzpicture}[->,>=stealth',shorten >=1pt,auto,node distance=2.8cm,
                    semithick]
  \tikzstyle{every state}=[fill=none,draw=black,text=black]

  \node[initial,state]         (A)                     {$q_0$};
  \node[state]         (B) [right of=A] {$q_1$};
  \node[state]         (C) [below of=A] {$h$};
  \node[state]         (D) [below of=B] {$r$};

  \path (A) edge [bend left] node {\TMtrans{\TMstroke}{\TMstroke}{R}} (B)
  	    edge node {\TMtrans{\TMblank}{\TMblank}{N}} (C)
        (B) edge node {\TMtrans{\TMblank}{\TMblank}{N}} (D)
            edge [bend left] node {\TMtrans{\TMstroke}{\TMstroke}{R}} (A);
\end{tikzpicture}
\]
\end{ex}

\begin{explain}
Nambah kahanan mandheg mligi bisa migunani
ing kahanan kaya iki, amarga kanthi mangkono
cetha kapan mesin nampa utawa nolak input
tartamtu. Nanging perlu digatekake yen
nambah kahanan mandheg mligi ora nambah
daya komputasi. Semono uga, pangerten
mandheg kang ora nganggo kahanan mligi nduwé paedahe
dhewe. Definisi mandheg kang wis digunakake
nganti saiki ing bab iki ndadekake bukti
\emph{Masalah Mandheg} gampang dimangerteni
lan dituduhake. Mula kita tetep nganggo
definisi asal.
\end{explain}

\end{document}
